\documentclass[12pt]{article}

% ---------- Packages ----------
\usepackage{fontspec} % Required for LuaLaTeX
\setmainfont{Times New Roman}

\usepackage[a4paper, margin=1in]{geometry}
\usepackage{setspace}
\usepackage{titlesec}
\usepackage{ragged2e}

% ---------- Formatting ----------
\setstretch{1.5} % Line spacing

% Section formatting (14pt, bold, uppercase)
\titleformat{\section}
  {\bfseries\fontsize{14}{16}\selectfont}
  {\thesection}{1em}{\MakeUppercase}

% Subsection formatting (12pt, bold, uppercase)
\titleformat{\subsection}
  {\bfseries\fontsize{12}{14}\selectfont}
  {\thesubsection}{1em}{\MakeUppercase}

% ---------- Document ----------
\begin{document}

% ---------- Title Page ----------
\begin{center}
    {\fontsize{16}{18}\selectfont \textbf{CS/AI/CY/SE4092}}\\[1.5em]

    {\fontsize{18}{20}\selectfont \textbf{Final Year Project}}\\[1em]

    {\fontsize{16}{18}\selectfont \textbf{Orbis AI}}\\[0.5em]
    {\fontsize{13}{15}\selectfont \textbf{FYP ID: SP-88}}\\[2em]

    {\fontsize{14}{16}\selectfont \textbf{Submitted By}}\\[1em]

    M. Talha Yousif (Leader) \hspace{1em} 22K-5146\\
    Salik Ahmed Bhatti \hspace{1em} 22K-5114\\
    Toheed Ali \hspace{1em} 22K-4027\\[2em]

    {\fontsize{16}{18}\selectfont \textbf{Project Progress Report}}\\[2em]

    Department of Computer Science\\
    National University of Computer and Emerging Sciences\\
    FAST Karachi Campus
\end{center}

\vspace{2em}

\begin{center}
\textbf{Project Group Members}
\end{center}

\vspace{0.5cm}

\begin{tabular}{|c|c|p{3.5cm}|c|p{3.5cm}|c|}
\hline
\textbf{Sr.\#} & \textbf{Reg. \#} & \textbf{Student Name} & \textbf{CGPA} & \textbf{Email ID} & \textbf{Phone \#} \\
\hline
(i) & 22K-5146 & M. Talha Yousif (Leader) & 2.76 & k225146@nu.edu.pk & 03164201069 \\[0.8cm]
\hline
(ii) & 22K-5114 & Salik Ahmed Bhatti & 2.26 & k5114@nu.edu.pk & 0322332833 \\[0.8cm]
\hline
(iii) & 22K-4027 & Toheed Ali & 2.16 & k224027@nu.edu.pk & 03043390057 \\[0.8cm]
\hline
\end{tabular}

\vspace{1.5em}

% ---------- Sections ----------

\section{Introduction}
\justifying
Orbis AI is an intelligent travel planning and booking assistant that combines conversational AI, retrieval-augmented generation (RAG), and task-specific agent workflows to help users plan end-to-end trips. The project objective is to move beyond simple chatbot interactions and provide practical travel support including itinerary planning, contextual recommendations, and booking-related actions through a unified interface.

The system is built around a FastAPI backend with LangGraph-based orchestration, Supabase for persistence, Redis caching, and a Next.js frontend for interactive chat. A major focus of this phase is strengthening system reliability and response quality by improving intent routing, adding conversation-aware context, and introducing structured agent tooling for flights, hotels, trips, and database operations.

Our recent improvements also include a significantly enhanced chat interface that supports real-time server-sent event (SSE) streaming, better message rendering, side-panel artifacts, and smoother user interactions. These upgrades make the product closer to a production-ready travel companion instead of a prototype.

\section{Timeline}
The high-level timeline from FYP Part 1 final to current progress is as follows:
\begin{itemize}
    \item \textbf{Phase 1 (Foundation):} Requirements finalization, architecture design, technology stack selection, and initial UI wireframes.
    \item \textbf{Phase 2 (Core Backend):} FastAPI setup, Supabase integration, authentication modules, conversation/message CRUD, and initial agent orchestration.
    \item \textbf{Phase 3 (Intelligence Layer):} LangGraph workflow improvements, intent routing expansion, RAG pipeline implementation, and embedding-based context retrieval.
    \item \textbf{Phase 4 (Product Refinement):} Chat interface improvements, streaming stability, additional routers (payments, analytics, notifications), and observability enhancements.
\end{itemize}

\section{Progress}
Current progress is mapped to milestone-driven development and shows strong advancement in backend capability, frontend usability, and overall system integration.

\subsection{Milestone 1 Progress}
Milestone 1 focused on core platform setup and baseline functionality. The following items were completed:
\begin{itemize}
    \item FastAPI server configured with CORS and SSE support.
    \item Supabase database connected with service-role access and data persistence.
    \item Authentication flows completed (register, login, refresh, profile).
    \item Conversation and message CRUD endpoints completed.
    \item Basic chat streaming endpoint created and integrated.
\end{itemize}

\subsection{Milestone 2 Progress}
Milestone 2 emphasized intelligence and orchestration quality:
\begin{itemize}
    \item LangGraph-based multi-node flow stabilized (orchestrator, planner, flight, hotel, itinerary, booking, verifier).
    \item Intent routing expanded to cover 20+ intent variants.
    \item Conversation history and date-awareness injected into system prompts.
    \item RAG pipeline integrated using embeddings and travel guide retrieval.
    \item Agent tooling definitions completed and integrated with ReAct-compatible workflow.
\end{itemize}

In this milestone, the system moved from static responses to context-aware travel assistance with significantly improved relevance and structured reasoning.

\subsection{Milestone 3 Progress}
Milestone 3 addressed product maturity, observability, and frontend experience:
\begin{itemize}
    \item Added new backend routers for payments, notifications, analytics, and agent feedback logging.
    \item Improved reliability through better conversation handling and cache behavior.
    \item Added automatic conversation title generation after early user-assistant exchanges.
    \item Enhanced chat interface with:
    \begin{itemize}
        \item smoother real-time streaming in the conversation view,
        \item improved message rendering/content parts handling,
        \item upgraded chat header and message view components,
        \item travel card and side-panel artifact presentation for richer UX.
    \end{itemize}
\end{itemize}

Overall, this milestone delivered significant frontend polish while preserving backend scalability and maintainability.

\section{Updated Timeline}
The updated timeline from proposal to expected completion is:
\begin{itemize}
    \item \textbf{Proposal \& Planning:} Problem definition, objectives, and architecture approval.
    \item \textbf{Implementation Sprint 1:} Backend/API foundation, auth, and database operations.
    \item \textbf{Implementation Sprint 2:} Agent orchestration, intent routing, and RAG integration.
    \item \textbf{Implementation Sprint 3:} Chat interface improvements, observability, and user experience refinement.
    \item \textbf{Remaining Work Before Final Completion:} Live external API credentials integration (Amadeus, hotel provider, Stripe), tool-call logging completion, timeout hardening, and comprehensive testing.
\end{itemize}

This updated timeline reflects a substantial completion level, with remaining tasks primarily focused on production readiness and validation.

\section*{References}
\begin{itemize}
    \item Orbis AI Internal Status Report, April 13, 2026.
    \item FastAPI Documentation. \textit{https://fastapi.tiangolo.com/}
    \item LangGraph and LangChain Documentation. \textit{https://python.langchain.com/}
    \item Supabase Documentation. \textit{https://supabase.com/docs}
    \item Next.js Documentation. \textit{https://nextjs.org/docs}
\end{itemize}

\end{document}